\section{Архитектура механизмов делегирования}

Правило имеет вид «шаблон $\Rightarrow$ вывод» с ограничениями-термами. 
В процессе применения его параметры получают значения: из шаблона (связанные) 
или из ограничений (свободные). Делегирование выражается в том, что 
ограничение или подвыражение вывода содержит вызов подзадачи либо внешнего 
алгоритма.

\subsection{Вызов подзадачи}

Подзадача --- вспомогательная задача, решаемая отдельной сессией решателя. 
Её результат подставляется в параметр правила. Подзадачи кэшируются; 
при повторном обращении к той же подзадаче возвращается маркер 
«решение не найдено», что обрывает рекурсию.

Правило может запускать подзадачу типа \texttt{find} --- поиск значений 
заданных переменных при заданных условиях. Контекст родительской задачи 
(известные факты) передаётся в подзадачу, но её цель независима. 
Рекурсия конечна, если подзадача строго проще исходной (например, 
понижение степени уравнения).

\subsection{Вызов внешнего алгоритма}

Внешний алгоритм оборачивается в \emph{процедурный символ} --- терм 
с произвольной арностью, тело которого вычисляется внешней функцией. 
Правило не различает процедурный и обычный символ. Процедурный символ 
в шаблоне не вычисляется, а участвует в унификации как константа; 
вычисление происходит только при его появлении в ограничении или выводе 
после связывания всех аргументов.

Если параметр в позиции вызова свободен, результат подставляется ему 
как значение. Если результат --- множество, оно раскрывается в ветви 
(декартово произведение). Если параметр уже связан или на месте вызова 
константа, ограничение проверяется на истинность; невыполненные ветви 
отбрасываются.

\subsection{Пример: правило рациональных корней}

Правило для монического многочлена с целыми коэффициентами:

\begin{lstlisting}
F == 0 => (x == u) || (Q == 0);
is_polynomial(F, x), d == free_term(F, x), d != 0,
u in divisors(d), substitute(x == u, F) == 0,
Q == poly_div(F, x - u)
\end{lstlisting}

Шаблон связывает $F$ с многочленом. \texttt{divisors(d)} возвращает 
множество делителей свободного члена --- порождает ветви. 
\texttt{substitute} проверяет обнуление; проходят только те $u$, 
при которых многочлен обращается в ноль. Вывод --- дизъюнкция: 
левая ветвь фиксирует корень, правая редуцирует задачу к многочлену 
меньшей степени.

\section{Гибридный подход при решении параметрических задач}

\subsection{Проблематика}

SymPy 1.14.0 систематически теряет решения в задачах с параметром. 
Типичный механизм потери --- деление на выражение, содержащее параметр, 
без отслеживания случая его обнуления. Например, 
\texttt{solve(a*x = a, x)} возвращает \texttt{[1]}, теряя случай 
$a = 0$.

Ниже используются три репрезентативные задачи (C1--C3), на которых 
SymPy даёт неполный ответ.

\subsection{Приёмы предобработки и постобработки}

\textbf{Предобработка} (до делегирования):

\begin{itemize}
  \item \emph{Факторизация}: выражение приводится к виду произведения. 
    Для уравнения $(a-2)(a+2)x = a-2$ разбор случаев по обнулению 
    $a-2$ и $a+2$ даёт полное решение.
  \item \emph{Разбор случаев}: вывод правила записывается как дизъюнкция. 
    Ветви порождаются либо связыванием параметра с множеством 
    (через $u \in \texttt{divisors}(d)$), либо явным перечислением.
\end{itemize}

\textbf{Постобработка} (после делегирования):

\begin{itemize}
  \item \emph{ОДЗ-фильтр}: кандидаты из ответа SymPy проверяются на 
    условия допустимости (например, знаменатель не ноль, подкоренное 
    выражение неотрицательно); неподходящие отбрасываются.
\end{itemize}

\subsection{Примеры}

\textbf{C1. Параметрическое линейное уравнение $a \cdot x = a$.}  
SymPy возвращает $\{1\}$, теряя $a = 0$. Гибридное правило:

\begin{lstlisting}
x*a == a =>
    answer(
        (a == 0 && x in Reals) ||
        (a != 0 && x in sympy_solve(x*a == a, x))
    )
\end{lstlisting}

Случаи перечислены явно. В ветви $a \neq 0$ \texttt{sympy\_solve} 
возвращает $\{1\}$. Итоговый ответ: 
$(a = 0 \wedge x \in \mathbb{R}) \vee (a \neq 0 \wedge x = 1)$.

\textbf{C2. $(a^2-4) \cdot x = a-2$.}  
SymPy возвращает $1/(a+2)$, теряя $a = 2$ (все действительные $x$) 
и $a = -2$ (нет решений). После факторизации $a^2-4 = (a-2)(a+2)$ 
применяется то же правило C1, давая полный кусочный ответ.

\textbf{C3. $(a^2-1)x^2 + 2(a-1)x + 1 = 0$.}  
SymPy возвращает формулу с $(a^2-1)$ в знаменателе, не разбирая случаи 
$a = \pm 1$. Гибридное правило:

\begin{lstlisting}
a*x^2 + b*x + c == 0 =>
    answer(
        (a == 0 && solution_linear) ||
        (a != 0 && x in sympy_solve(a*x^2 + b*x + c == 0, x))
    );
a, b, c known,
solution_linear == find(x) { b*x + c == 0 }
\end{lstlisting}

Подзадача \texttt{find(x) \{ b*x + c == 0 \}} решается правилом C1. 
Её ответ --- множество решений --- связывается с параметром 
\texttt{solution\_linear}. При $a = 0$ ответом служит результат 
подзадачи; при $a \neq 0$ делегирование идёт в SymPy.

\subsection{Результаты}

Гибридная система возвращает корректный кусочный ответ на C1--C3. 
На C1 и C2 разбор случаев по обнулению коэффициента устраняет потери SymPy. 
На C3 подзадача обрабатывает вырожденную ветвь ($a=0$), а CAS --- 
невырожденную ($a \neq 0$). SymPy 1.14.0 на всех трёх задачах теряет 
решения.